\section{Conclusion}
\label{sec:conclusion}

We tested whether check-in-level contextual graph embeddings carry the multi-task POI prediction win on next-category and next-region across five U.S.-state Gowalla splits. The substrate is task-asymmetric: per-visit context lifts matched-head single-task next-category macro-F1 by $+14.5$ to $+29$ pp at every state (paired Wilcoxon $p = 0.0312$, the $n = 5$ maximum); on next-region the same substrate ties or marginally trails per-place HGI by $1.6$ to $3.1$ pp under the matched-head \texttt{next\_stan\_flow} ceiling, with TOST non-inferiority at $\delta = 2$ pp passing at CA and TX. A pooled-vs-canonical counterfactual at Alabama attributes about $72\,\%$ of the category substrate gap to per-visit context, single-state mechanism evidence consistent with the asymmetry. With Check2HGI fixed, the cross-attention multi-task backbone adds a small category lift at four of five states (AZ $+1.20$, FL $+1.40$, CA $+1.68$, TX $+1.89$ pp at $p \leq 2\mathrm{e}{-6}$ across $n = 20$ fold-pairs) and is small-significantly negative at AL ($\Delta = -0.78$ pp, $p = 0.036$, $n = 20$); on next-region it pays a sign-consistent $7$ to $17$ pp cost on Acc@10 at every state ($p \leq 1.9\mathrm{e}{-6}$, $n = 20$), the textbook multi-task tradeoff. On the joint $\Delta_m$ axis, FL is the lone Pareto-positive cell on the primary (category macro-F1, region MRR) pair at $+2.33\,\%$ ($p = 2.98\mathrm{e}{-8}$, $25/25$ fold-pairs positive); the other four states sit below zero on both axes at the $n = 5$ ceiling. Drop-in alternatives (FAMO, Aligned-MTL, a hierarchical-additive-softmax region head) do not close the next-region gap, and the methodological note of Section~\ref{subsec:mechanism-isolation} generalises beyond our backbone: encoder-frozen isolation is the clean architectural decomposition under cross-attention.

Three forward-looking lines suggest themselves. Asymmetric routing through PLE~\cite{tang2020ple} or Cross-Stitch Networks~\cite{misra2016crossstitch} could free the region head from shared-backbone interference. Encoder enrichment, fusing per-visit context with explicit spatial and temporal sub-encoders, plausibly reaches a substrate that the joint head can read without sacrificing either task. A region-side dynamic prior, conditioned on the user's recent trajectory, would target the next-region cost at the head rather than at the substrate. Recovering the next-region performance without sacrificing the next-category lift is the open challenge this paper leaves to follow-up work.
